\section{Особенности реализации}

Основные алгоритмы реализованы методами класса \texttt{GeneratorGraph}. Взаимодействие с пользователем организовано классом \texttt{CLI} через общее консольное меню.

Далее через $p$ обозначается количество вершин графа, через $q$~--- количество его рёбер, а через $k$~--- заданное количество рёбер маршрута.

\subsection{Лабораторная работа № 1}

Алгоритм генерации графа реализован методом \texttt{generate}, приведённым в листинге~\ref{lst:generate-graph}. Он формирует связный ориентированный граф по последовательности степеней, полученной с помощью распределения Вейбулла.

\textbf{Вход:} количество вершин графа.

\textbf{Выход:} матрицы смежности ориентированного и неориентированного графов.

\textbf{Сложность:} $O(p^2)$.

\codeListing[cpp][firstline=85,lastline=118,breaklines]{../src/generator_graph.cpp}{Генерация графа}{\label{lst:generate-graph}}

Алгоритм вычисления эксцентриситетов реализован методом \texttt{computeEccentricities}, приведённым в листинге~\ref{lst:compute-eccentricities}. По степеням матрицы смежности он находит длины кратчайших маршрутов и вычисляет характеристики вершин.

\textbf{Вход:} матрица смежности графа.

\textbf{Выход:} эксцентриситеты, диаметр, центральные и диаметральные вершины.

\textbf{Сложность:} $O(p^4)$, поскольку выполняется до $p-1$ умножений матриц.

\codeListing[cpp][firstline=120,lastline=159,breaklines]{../src/generator_graph.cpp}{Вычисление эксцентриситетов, центра и диаметра}{\label{lst:compute-eccentricities}}

Метод \texttt{generateWeight\allowbreak{}Matrix}, приведённый в листинге~\ref{lst:generate-weight-matrix}, присваивает существующим рёбрам веса в выбранном режиме.

\textbf{Вход:} матрица смежности и режим генерации положительных, отрицательных или смешанных весов.

\textbf{Выход:} весовые матрицы ориентированного и неориентированного графов.

\codeListing[cpp][firstline=166,lastline=202,breaklines]{../src/generator_graph.cpp}{Генерация весовой матрицы}{\label{lst:generate-weight-matrix}}

Метод Шимбелла реализован методом \texttt{shimbell}, приведённым в листинге~\ref{lst:shimbell-implementation}. Он последовательно строит весовую матрицу маршрутов заданной длины.

\textbf{Вход:} матрицы смежности и весов, количество рёбер маршрута и требуемый тип веса: минимальный либо максимальный.

\textbf{Выход:} матрица весов маршрутов заданной длины.

\textbf{Сложность:} $O(kp^3)$.

\codeListing[cpp][firstline=204,lastline=237,breaklines]{../src/generator_graph.cpp}{Реализация метода Шимбелла}{\label{lst:shimbell-implementation}}

Алгоритм исследования маршрутов реализован методами \texttt{hasRoute} и \texttt{countRoutes}, приведёнными в листинге~\ref{lst:routes-implementation}. Методы используют степени матрицы смежности для исследования маршрутов между двумя вершинами.

\textbf{Вход:} матрица смежности, начальная и конечная вершины.

\textbf{Выход:} признак существования маршрута и количество найденных маршрутов.

\textbf{Сложность:} $O(p^4)$.

\codeListing[cpp][firstline=239,lastline=260,breaklines]{../src/generator_graph.cpp}{Проверка наличия и подсчёт маршрутов}{\label{lst:routes-implementation}}

\subsection{Лабораторная работа № 2}

Алгоритм поиска точек сочленения реализован методами \texttt{m\_dfsArticulation} и \texttt{findArticulationPoints}, приведёнными в листинге~\ref{lst:articulation-implementation}. Они выполняют обход графа в глубину, вычисляют время входа и нижние метки вершин.

\textbf{Вход:} матрица смежности.

\textbf{Выход:} список точек сочленения и количество итераций алгоритма.

\textbf{Сложность:} $O(p^2)$ при хранении графа в виде матрицы смежности.

\codeListing[cpp][firstline=262,lastline=333,breaklines]{../src/generator_graph.cpp}{Поиск точек сочленения}{\label{lst:articulation-implementation}}

Модифицированный алгоритм Дейкстры реализован методом \texttt{dijkstraNegative}, приведённым в листинге~\ref{lst:dijkstra-negative-implementation}. Он выбирает из очереди вершину с минимальным текущим весом, допускает её повторную обработку и контролирует число уменьшений веса для обнаружения отрицательного цикла.

\textbf{Вход:} матрица смежности, весовая матрица, начальная и конечная вершины.

\textbf{Выход:} вектор расстояний, предшественники вершин, кратчайший путь, его длина, число итераций и признак отрицательного цикла.

\textbf{Сложность:} $O(e^p)$ в худшем случае.

\codeListing[cpp][firstline=369,lastline=458,breaklines]{../src/generator_graph.cpp}{Поиск кратчайшего пути при наличии отрицательных весов}{\label{lst:dijkstra-negative-implementation}}

\subsection{Лабораторная работа № 3}

Метод \texttt{generateCapacityAndCostMatrices}, приведённый в листинге~\ref{lst:capacity-cost-generation}, назначает каждому ребру исходного ориентированного графа пропускную способность и стоимость.

\textbf{Вход:} матрица смежности ориентированного графа.

\textbf{Выход:} матрицы пропускных способностей и стоимостей.

\codeListing[cpp][firstline=480,lastline=501,breaklines]{../src/generator_graph.cpp}{Генерация пропускных способностей и стоимостей}{\label{lst:capacity-cost-generation}}

Алгоритм Форда~--- Фалкерсона реализован методом \texttt{fordFulkerson}, приведённым в листинге~\ref{lst:ford-fulkerson-implementation}. Он расставляет метки, находит увеличивающую цепь и изменяет поток вдоль неё, пока сток может получить метку.

\textbf{Вход:} матрица пропускных способностей, исток и сток.

\textbf{Выход:} величина максимального потока, матрица потоков и количество итераций.

\textbf{Сложность:} $O(Ap^3)$, где $A$~--- количество найденных увеличивающих цепей.

\codeListing[cpp][firstline=531,lastline=675,breaklines]{../src/generator_graph.cpp}{Вычисление максимального потока}{\label{lst:ford-fulkerson-implementation}}

Алгоритм последовательного нахождения потока минимальной стоимости реализован методом \texttt{minCostFlow}, приведённым в листинге~\ref{lst:min-cost-flow-implementation}. Он находит кратчайшие пути в сети и изменяет поток до достижения требуемой величины либо отсутствия подходящего пути.

\textbf{Вход:} матрицы пропускных способностей и стоимостей, исток, сток и требуемая величина потока.

\textbf{Выход:} достигнутая величина и стоимость потока, матрицы потоков и стоимостей потока, число итераций и признак успешного завершения.

\textbf{Сложность:} $O(Ae^p)$ в худшем случае, где $A$~--- количество увеличений потока. Экспоненциальная сложность обусловлена использованием модифицированного алгоритма Дейкстры при каждом увеличении потока.

\codeListing[cpp][firstline=677,lastline=779,breaklines]{../src/generator_graph.cpp}{Вычисление потока минимальной стоимости}{\label{lst:min-cost-flow-implementation}}

\subsection{Лабораторная работа № 4}

Алгоритм подсчёта остовных деревьев по теореме Кирхгофа реализован методом \texttt{countSpanningTreesKirchhoff}, приведённым в листинге~\ref{lst:kirchhoff-implementation}. Он строит матрицу Лапласа и вычисляет определитель её минора.

\textbf{Вход:} матрица смежности неориентированного графа.

\textbf{Выход:} матрица Лапласа, её минор и число остовных деревьев.

\textbf{Сложность:} $O(p!)$, поскольку определитель вычисляется рекурсивным разложением по строке.

\codeListing[cpp][firstline=782,lastline=819,breaklines]{../src/generator_graph.cpp}{Подсчёт остовных деревьев по теореме Кирхгофа}{\label{lst:kirchhoff-implementation}}

Определитель вычисляется методом \texttt{determinant}, приведённым в листинге~\ref{lst:determinant-method}. Метод проверяет, что матрица является квадратной, и вызывает рекурсивное разложение, приведённое в листинге~\ref{lst:determinant-recursive}. Для сокращения количества слагаемых выбирается строка или столбец с наибольшим числом нулевых элементов.

\textbf{Вход:} квадратная матрица.

\textbf{Выход:} определитель матрицы.

\textbf{Сложность:} $O(p!)$ в худшем случае.

\codeListing[cpp][firstline=192,lastline=201,breaklines,fontsize=\footnotesize]{../src/matrix.cpp}{Метод вычисления определителя}{\label{lst:determinant-method}}

\codeListing[cpp][firstline=13,lastline=133,breaklines,fontsize=\footnotesize]{../src/matrix.cpp}{Рекурсивное разложение определителя}{\label{lst:determinant-recursive}}

Алгоритм Краскала реализован методами \texttt{m\_kruskalAlgorithm} и \texttt{kruskalMinimumSpanningTree}, приведёнными в листинге~\ref{lst:kruskal-implementation}. Они сортируют рёбра, объединяют множества вершин и формируют минимальный остов.

\textbf{Вход:} матрицы смежности и весов неориентированного графа.

\textbf{Выход:} рёбра и вес минимального остова, его весовая матрица и признак успешного построения.

\textbf{Сложность:} $O(p^2+q\log q)$ с учётом формирования списка рёбер из матрицы смежности.

\codeListing[cpp][firstline=821,lastline=937,breaklines]{../src/generator_graph.cpp}{Построение минимального остова алгоритмом Краскала}{\label{lst:kruskal-implementation}}

Алгоритмы кодирования и декодирования кода Прюфера реализованы методами \texttt{m\_encodePruferCode} и \texttt{m\_decodePruferCode}, приведёнными в листинге~\ref{lst:prufer-implementation}. Они кодируют минимальный остов и восстанавливают его с сохранением весов рёбер.

\textbf{Вход:} минимальный остов при кодировании или код Прюфера при декодировании.

\textbf{Выход:} код Прюфера либо восстановленная весовая матрица остова.

\textbf{Сложность:} $O(p^2)$ для кодирования и $O(p^2\log p)$ для декодирования.

\codeListing[cpp][firstline=939,lastline=1031,breaklines]{../src/generator_graph.cpp}{Кодирование и декодирование остова}{\label{lst:prufer-implementation}}

Алгоритм Эдмондса реализован методами \texttt{edmondsBlossom\allowbreak{}InOriginalGraph}, \texttt{edmondsBlossom\allowbreak{}InMinimumSpanningTree} и \texttt{m\_edmondsBlossom}, приведёнными в листинге~\ref{lst:edmonds-implementation}. Они выбирают исследуемый граф и находят максимальное независимое множество рёбер, сжимая обнаруженные нечётные циклы.

\textbf{Вход:} матрицы смежности и весов исходного графа либо минимального остова.

\textbf{Выход:} список и матрица рёбер максимального паросочетания.

\textbf{Сложность:} $O(p^3)$.

\codeListing[cpp][firstline=1033,lastline=1267,breaklines]{../src/generator_graph.cpp}{Поиск максимального независимого множества рёбер}{\label{lst:edmonds-implementation}}

\subsection{Лабораторная работа № 5}

Алгоритм приведения графа к эйлерову виду реализован методом \texttt{checkIfEulerianGraph}, приведённым в листинге~\ref{lst:eulerian-transformation-implementation}. Он проверяет чётность степеней и попарно добавляет либо удаляет рёбра между нечётными вершинами. Перед удалением проверяется наличие обходного пути, поэтому преобразование прекращается, если оно требует удаления моста.

\textbf{Вход:} матрица смежности неориентированного графа.

\textbf{Выход:} изменённая матрица смежности, списки добавленных и удалённых рёбер, признаки исходного состояния и завершения преобразования.

\textbf{Сложность:} в худшем случае $O(p^6)$ с учётом перебора пар нечётных вершин и повторных поисков обходного пути.

\codeListing[cpp][firstline=1281,lastline=1373,breaklines]{../src/generator_graph.cpp}{Проверка и преобразование графа к эйлерову виду}{\label{lst:eulerian-transformation-implementation}}

Алгоритм Флёри реализован методом \texttt{m\_flueryAlgorithm}, приведённым в листинге~\ref{lst:fleury-implementation}. Он последовательно выбирает соседнюю вершину, удаляет пройденное ребро и добавляет вершину в результат после исчерпания её соседей.

\textbf{Вход:} матрица смежности эйлерова графа.

\textbf{Выход:} последовательность вершин эйлерова цикла.

\textbf{Сложность:} $O(pq)$ в худшем случае.

\codeListing[cpp][firstline=1375,lastline=1417,breaklines]{../src/generator_graph.cpp}{Построение эйлерова цикла}{\label{lst:fleury-implementation}}

Алгоритм построения фундаментальной системы разрезов реализован методом \texttt{buildFundamentalCutSystem}, приведённым в листинге~\ref{lst:fundamental-cuts-implementation}. Он строит минимальный остов и для каждого его ребра получает фундаментальный разрез.

\textbf{Вход:} матрицы смежности и весов неориентированного графа.

\textbf{Выход:} минимальный остов и фундаментальная система разрезов.

\textbf{Сложность:} $O(p^3+q\log q)$.

\codeListing[cpp][firstline=1419,lastline=1482,breaklines]{../src/generator_graph.cpp}{Построение фундаментальной системы разрезов}{\label{lst:fundamental-cuts-implementation}}

Алгоритм получения остальных разрезов реализован методами \texttt{symmetric\allowbreak{}Difference\allowbreak{}Of\allowbreak{}Fundamental\allowbreak{}Cuts} и \texttt{m\_symmetric\allowbreak{}Difference}, приведёнными в листинге~\ref{lst:cut-symmetric-difference-implementation}. Они последовательно вычисляют симметрическую разность выбранных фундаментальных разрезов.

\textbf{Вход:} фундаментальная система разрезов и выбранные разрезы.

\textbf{Выход:} множество рёбер результирующего разреза.

\textbf{Сложность:} $O(rq\log q)$, где $r$~--- количество выбранных фундаментальных разрезов.

\codeListing[cpp][firstline=1484,lastline=1528,breaklines]{../src/generator_graph.cpp}{Получение разреза симметрической разностью}{\label{lst:cut-symmetric-difference-implementation}}
